\section{Kilka słów wstępu}
\textit{Labirynt śmierci} jest grą nowej generacji, zupełnie niepodobną do jakiejkolwiek innej znanej na polskim rynku. Wymaga od grających dużego zaangażowania intelektualnego i skupienia uwagi podczas całej rozgrywki~--- w zamian dostarcza niezapomnianych emocji. Kryje w sobie wiele niespodzianek, które powodują, że gra się nie nudzi~--- a wręcz przeciwnie~--- staje się coraz bardziej atrakcyjna w miarę jak gracze nabierają doświadczenia. Eksplorację labiryntu, poszukiwanie kosztowności i magicznych przedmiotów, walkę z potworami, rzucanie i odpieranie czarów można bowiem realizować w różny sposób; zaś efekty są zależne od przyjętego sposobu postępowania.

Lecz \textit{Labirynt śmierci} jest grą nietypową przede wszystkim dlatego, że gracze (od 1 do 6 osób) nie grają przeciwko sobie. Wręcz przeciwnie~--- aby osiągnąć cel, tzn. odnaleźć i zniszczyć Czarne Wrota, musza wzajemnie się wspierać i ratować w trudnych momentach. Gra toczy się przeciwko sytuacji, w jakiej gracze, czy może raczej występujące w grze postacie się znajdą, przeciwko niebezpieczeństwom, potworom, pułapkom, magii. Tylko ścisła wspólpraca wszystkich grających może zapewnić odniesienie zwycięstwa. Gracz nastawiony wyłącznie na idywidualny sukces nie tylko nie osiągnie celu, ale może spowodować śmierć postaci, którą kieruje.

Instrukcja, jak nietrudno zauważyć, jest długa. Dlatego pozwolimy sobie na jedną praktyczną radę~--- prosimy nie starać się opanować wszystkich zawartych w niej zasad od razu, przed rozpoczęciem gry. Znacznie bardziej efektywna metoda ich nauczenia się polega na jedno- lub dwukrotnym przeczytaniu instrukcji, by wyrobić sobie ogólne pojęcie o grze, rządzących nią prawach i jej przebiegu~--- a następnie na przystąpieniu do rozgrywki i w chwili, gdy pojawią się niejasności~--- na konsultowaniu się z odpowiednim paragrafem instrukcji.
